\chapter{CƠ SỞ LÝ THUYẾT}
\label{chap:background}

\section{Mô hình ngôn ngữ tự hồi quy}
\label{sec:bg-lm}

Một mô hình ngôn ngữ (LM) gán xác suất cho chuỗi token $x_{1:T}$ bằng phân rã
tự hồi quy:
\begin{equation}
	p_\theta(x_{1:T}) = \prod_{t=1}^{T} p_\theta(x_t \mid x_{<t}),
	\label{eq:autoregressive}
\end{equation}
tức toàn bộ việc ``hiểu ngôn ngữ'' được quy về một bài toán duy nhất:
\emph{dự đoán token kế tiếp}. Huấn luyện tối thiểu hóa cross-entropy trên
từng vị trí:
\begin{equation}
	\mathcal{L}(\theta) = -\frac{1}{T}\sum_{t=1}^{T} \log p_\theta(x_t \mid x_{<t}).
	\label{eq:ce}
\end{equation}
Điểm đáng chú ý về mặt sư phạm: công thức~\eqref{eq:ce} \emph{trùng khớp} với
loss của decoder sinh mô tả ảnh ở đồ án cuối kỳ --- chỉ khác là ở đó điều kiện
$x_{<t}$ được nối thêm đặc trưng ảnh. Mô hình ngôn ngữ vì vậy có thể xem là
``captioning bỏ điều kiện ảnh, phóng đại dữ liệu lên nghìn lần''; và ngược
lại, VLM ở chương~\ref{chap:methods} chính là mô hình ngôn ngữ \emph{thêm lại}
điều kiện ảnh.

Hai thước đo chất lượng nội tại được dùng trong báo cáo. \textbf{Perplexity}
$\mathrm{PPL} = \exp(\mathcal{L})$ đo mức ``bối rối'' trung bình trên mỗi
token; tuy nhiên PPL phụ thuộc bộ tách từ: cùng một đoạn văn, tokenizer cắt
càng thô (ít token hơn) thì mỗi token càng khó đoán, PPL càng cao một cách tự
nhiên, nên \emph{không thể so PPL giữa hai mô hình dùng hai tokenizer khác
nhau}. Vì vậy khi so với baseline, báo cáo dùng \textbf{bits-per-character}:
\begin{equation}
	\mathrm{BPC} = \frac{\sum_t -\log_2 p_\theta(x_t \mid x_{<t})}{\text{số ký tự của văn bản}},
	\label{eq:bpc}
\end{equation}
tổng lượng thông tin (bit) mô hình cần để nén văn bản, chia cho độ dài
\emph{ký tự} --- đại lượng trung lập với mọi cách tách token.

\section{Kiến trúc decoder-only hiện đại}
\label{sec:bg-arch}

\subsection{Khối transformer và attention nhân quả}

Mô hình trong báo cáo là transformer decoder-only
\autocite{vaswani2017attention,radford2019gpt2}: $L$ khối xếp chồng, mỗi khối
gồm self-attention đa đầu và mạng truyền thẳng (FFN), nối tắt phần dư và chuẩn
hóa pre-norm. Self-attention tính
$\mathrm{softmax}(QK^\top/\sqrt{d_k})V$ với mặt nạ \emph{nhân quả} chặn mọi vị
trí tương lai, bảo đảm tính tự hồi quy của~\eqref{eq:autoregressive}. Cài đặt
dùng \texttt{scaled\_dot\_product\_attention} của PyTorch, tự kích hoạt nhân
FlashAttention \autocite{dao2022flashattention} --- tính attention chính xác
nhưng không ghi ma trận $T \times T$ ra bộ nhớ, yếu tố quyết định để đẩy
kích thước batch lên cao trong 24\,GB VRAM.

\subsection{Ba nâng cấp Llama-style so với GPT-2}

So với GPT-2 nguyên bản, kiến trúc trong báo cáo theo ``chuẩn thực hành''
hiện nay (LLaMA \autocite{touvron2023llama} và hầu hết mô hình mở sau nó),
khác ở ba điểm, tóm tắt trong bảng~\ref{tab:gpt2-vs-llama}.

\begin{table}[H]
	\centering
	\caption{Ba khác biệt kiến trúc giữa GPT-2 nguyên bản và Llama-style.}
	\label{tab:gpt2-vs-llama}
	\begin{tabularx}{\textwidth}{lXX}
		\hline
		\textbf{Thành phần} & \textbf{GPT-2} & \textbf{Llama-style (báo cáo này)} \\
		\hline
		Mã hóa vị trí & Embedding vị trí học được, cộng vào đầu vào & RoPE: xoay từng cặp chiều của $Q,K$ một góc tỷ lệ vị trí \\
		Chuẩn hóa & LayerNorm (trừ trung bình, có bias) & RMSNorm: chỉ chia RMS, không trừ trung bình, không bias \\
		FFN & Hai lớp tuyến tính với GELU & SwiGLU: cổng nhân điểm $\mathrm{SiLU}(xW_1)\odot xW_3$ rồi chiếu xuống \\
		\hline
	\end{tabularx}
\end{table}

\textbf{RoPE} \autocite{su2024rope} chia mỗi vector query/key thành các cặp
chiều $(2i, 2i{+}1)$ và xoay cặp thứ $i$ tại vị trí $t$ một góc
$t\,\omega_i$:
\begin{equation}
	\begin{pmatrix} q'_{2i} \\ q'_{2i+1} \end{pmatrix}
	= \begin{pmatrix} \cos t\omega_i & -\sin t\omega_i \\ \sin t\omega_i & \cos t\omega_i \end{pmatrix}
	  \begin{pmatrix} q_{2i} \\ q_{2i+1} \end{pmatrix},
	\qquad \omega_i = \theta^{-2i/d},\ \theta = 10^4.
	\label{eq:rope}
\end{equation}
Vì tích vô hướng giữa hai vector đã xoay chỉ phụ thuộc \emph{hiệu} hai góc,
attention nhìn thấy vị trí \emph{tương đối} $t - s$ thay vì tuyệt đối --- một
tiên nghiệm đúng với ngôn ngữ hơn embedding vị trí học được, lại không tốn
tham số. Phép xoay bảo toàn chuẩn vector (một bất biến được unit test khóa
trong mã nguồn).

\textbf{RMSNorm} \autocite{zhang2019rmsnorm} bỏ bước trừ trung bình của
LayerNorm, chỉ giữ $x / \sqrt{\tfrac{1}{d}\sum_j x_j^2 + \epsilon} \cdot g$:
rẻ hơn, ít tham số hơn (không bias), thực nghiệm cho chất lượng tương đương.
\textbf{SwiGLU} \autocite{shazeer2020glu} thay FFN hai lớp bằng ba ma trận với
cổng nhân: phần cổng $\mathrm{SiLU}(xW_1)$ điều tiết theo từng chiều lượng
thông tin đi qua $xW_3$, cho chất lượng tốt hơn trên cùng ngân sách FLOP; để
giữ tổng tham số tương đương, chiều ẩn giảm từ $4d$ xuống $\approx \tfrac{8}{3}d$.

Cuối cùng, \textbf{trói trọng số} (weight tying): ma trận embedding
$E \in \mathbb{R}^{|V| \times d}$ được dùng lại làm lớp chiếu đầu ra
($\mathrm{logits} = hE^\top$), tiết kiệm $|V| \times d$ tham số --- đáng kể ở
quy mô nhỏ, nơi embedding chiếm tỷ trọng lớn.

\section{Tách từ BPE và độ phì (fertility)}
\label{sec:bg-bpe}

Byte-Pair Encoding \autocite{sennrich2016bpe} khởi đầu từ bảng chữ cái và lặp
việc gộp cặp ký hiệu liền kề xuất hiện nhiều nhất thành mục từ vựng mới, đến
khi đạt kích thước từ vựng định trước. Biến thể \emph{byte-level} (GPT-2 và
báo cáo này) làm việc trên byte UTF-8 thay vì ký tự: mọi ký tự --- kể cả chữ
có dấu tiếng Việt như ``ắ'', ``ễ'' --- luôn phân rã được về byte, nên không
bao giờ xuất hiện token \texttt{<unk>}.

Hiệu quả của tokenizer đối với một ngôn ngữ được đo bằng \textbf{fertility}:
số token trung bình cần để mã hóa một từ. Fertility thấp nghĩa là chuỗi ngắn
hơn --- mô hình ``nhìn'' được ngữ cảnh dài hơn trong cùng cửa sổ, sinh nhanh
hơn, mỗi token mang nhiều nghĩa hơn. Một tokenizer huấn luyện trên tiếng Anh
(GPT-2) sẽ băm nát từ tiếng Việt có dấu thành nhiều mảnh byte; đây là lý do
báo cáo tự huấn luyện BPE 20.480 mục trên chính corpus tiếng Việt --- cùng
kích thước từ vựng với PhoGPT \autocite{nguyen2023phogpt} để so sánh trực
tiếp. Từ vựng nhỏ (20K thay vì 50K của GPT-2) còn có lợi ở quy mô nano:
embedding chỉ chiếm 15,7M thay vì 38M tham số, dành dung lượng cho các tầng
transformer.

\section{Quy luật scaling và điểm tối ưu tính toán}
\label{sec:bg-scaling}

Loss tiền huấn luyện giảm theo quy luật lũy thừa theo cả kích thước mô hình
lẫn lượng dữ liệu \autocite{kaplan2020scaling}. Với ngân sách tính toán cố
định, Chinchilla \autocite{hoffmann2022chinchilla} chỉ ra tỷ lệ tối ưu xấp xỉ
\emph{20 token dữ liệu cho mỗi tham số}. Mô hình 101M tham số của báo cáo do
đó ứng với điểm tối ưu $\approx 2$ tỷ token; lượng dữ liệu thực dùng (3,02 tỷ
token, $\approx 30$ token/tham số) chủ đích nằm \emph{quá} điểm Chinchilla ---
hợp lý khi mô hình (chứ không phải ngân sách huấn luyện) là thứ được giữ lại
dùng tiếp cho các giai đoạn sau, tương tự xu hướng ``over-train'' các mô hình
nhỏ thực dụng.

\section{Hợp nhất thị giác kiểu LLaVA}
\label{sec:bg-fusion}

\subsection{Bộ mã hóa thị giác tiền huấn luyện tương phản}

CLIP \autocite{radford2021clip} huấn luyện song song một bộ mã hóa ảnh và một
bộ mã hóa văn bản sao cho ảnh và mô tả đúng của nó có biểu diễn gần nhau
(học tương phản); kết quả là đặc trưng ảnh ``đã ngấm'' ngữ nghĩa ngôn ngữ ---
nền tảng cho hầu hết VLM. SigLIP \autocite{zhai2023siglip} thay softmax toàn
batch bằng loss sigmoid từng cặp, huấn luyện ổn định hơn; biến thể SigLIP-B/16
dùng trong báo cáo là một Vision Transformer \autocite{dosovitskiy2021vit}
chia ảnh $224 \times 224$ thành $14 \times 14 = 196$ patch, mỗi patch một
vector 768 chiều.

\subsection{Ghép token ảnh vào mô hình ngôn ngữ}

LLaVA \autocite{liu2023llava} chứng minh một cách hợp nhất tối giản mà hiệu
quả: chiếu các patch ảnh qua một MLP nhỏ (projector) vào đúng không gian
embedding của LLM, rồi \emph{ghép trước} (prepend) chuỗi token ảnh vào chuỗi
token văn bản. Từ đó self-attention của LLM tự nhiên attend chéo giữa từ và
vùng ảnh --- không cần khối cross-attention riêng như kiến trúc
encoder--decoder của đồ án cuối kỳ. Khối ``fusion'' trong sơ đồ của môn học,
theo cách nhìn này, chính là các tầng transformer sẵn có của mô hình ngôn ngữ.

Chi phí attention tỷ lệ bình phương độ dài chuỗi, nên số token ảnh đáng được
tiết kiệm. Kỹ thuật \emph{pixel-shuffle} (InternVL \autocite{chen2024internvl},
nanoVLM \autocite{wiedmann2025nanovlm}) gộp mỗi khối $2 \times 2$ patch lân
cận thành một token với số chiều gấp 4 ($196 \to 49$ token, $768 \to 3072$
chiều), giữ nguyên thông tin nhưng giảm 4 lần chi phí attention cho phần ảnh;
projector sau đó chiếu 3072 chiều về đúng chiều $d$ của mô hình ngôn ngữ.

\section{Tinh chỉnh có giám sát và che loss}
\label{sec:bg-sft}

Mô hình sau tiền huấn luyện chỉ biết \emph{tiếp tục} văn bản, chưa biết
\emph{trả lời}. SFT dạy định dạng hội thoại bằng chính loss
next-token~\eqref{eq:ce} trên dữ liệu có cấu trúc
``\texttt{<|user|>} câu hỏi \texttt{<|assistant|>} câu trả lời'', với một
điểm mấu chốt: \textbf{che loss} (loss masking) --- các vị trí thuộc phần
prompt (và token ảnh) bị gán nhãn bỏ qua, gradient chỉ chảy từ phần trả lời.
Mô hình vì vậy học ``khi thấy câu hỏi thì sinh câu trả lời'' chứ không học
thuộc lòng câu hỏi. Đây là khái niệm mà đồ án cuối kỳ chưa chạm tới (ở đó
toàn bộ caption là mục tiêu học), và là mảnh ghép ``SFT'' trong sơ đồ của
môn học.

\section{Học tăng cường cho sinh chuỗi: SCST}
\label{sec:bg-scst}

Các thước đo sinh chuỗi (CIDEr \autocite{vedantam2015cider}, BLEU
\autocite{papineni2002bleu}) không khả vi, nên không đưa thẳng vào~\eqref{eq:ce}
được. REINFORCE \autocite{williams1992reinforce} cho phép tối ưu kỳ vọng phần
thưởng $r$ của chuỗi lấy mẫu $w^s \sim p_\theta$:
\begin{equation}
	\nabla_\theta \mathbb{E}[r] \approx (r(w^s) - b)\,\nabla_\theta \log p_\theta(w^s),
	\label{eq:reinforce}
\end{equation}
với baseline $b$ để giảm phương sai. SCST \autocite{rennie2017scst} chọn
baseline là phần thưởng của chính mô hình khi giải mã tham lam $\hat{w}$:
$b = r(\hat{w})$ --- mẫu ngẫu nhiên nào tốt hơn ``chính mình lúc nghiêm túc''
thì được tăng xác suất, kém hơn thì bị giảm. Đồ án cuối kỳ đã kiểm chứng SCST
trên captioning tiếng Anh; báo cáo này tái sử dụng đúng thuật toán nhưng áp
lên VLM tự pretrain, khép mảnh ``RL'' trong sơ đồ.

\section{Công trình liên quan và vị trí của đề tài}
\label{sec:bg-related}

\textbf{Huấn luyện GPT tối giản.} nanoGPT \autocite{karpathy2022nanogpt} là
tham chiếu chuẩn cho tiền huấn luyện GPT quy mô nhỏ (kiến trúc, siêu tham số,
memmap dữ liệu); báo cáo kế thừa tinh thần tối giản này nhưng chuyển sang
kiến trúc Llama-style và tiếng Việt. Với tiếng Việt,
NlpHUST/gpt2-vietnamese \autocite{nlphust2021gpt2vi} (GPT-2 huấn luyện trên
tin tức) được dùng làm baseline BPC, còn PhoGPT-4B \autocite{nguyen2023phogpt}
--- tiền huấn luyện từ đầu 3,7B tham số trên cụm nhiều GPU --- là mốc trên
tham khảo và nguồn đối chứng thiết kế tokenizer.

\textbf{VLM nhỏ và VLM tiếng Việt.} nanoVLM \autocite{wiedmann2025nanovlm}
(SigLIP-85M + SmolLM2-135M, 222M tham số) là tiền lệ gần nhất về kiến trúc và
quy mô, nhưng \emph{mượn} LLM tiếng Anh huấn luyện sẵn. LaVy
\autocite{tran2024lavy} và Vintern-1B \autocite{doan2024vintern} là các VLM
tiếng Việt công bố sớm, đều xây trên LLM đa ngôn ngữ tỷ tham số. Bảng
\ref{tab:related} định vị ViVLM-nano: khoảng trống mà đề tài nhắm tới là tổ
hợp \emph{tự pretrain phần ngôn ngữ} + \emph{tiếng Việt} + \emph{trọn pipeline
ba giai đoạn} + \emph{một GPU}.

\begin{table}[H]
	\centering
	\caption{ViVLM-nano so với các hệ liên quan.}
	\label{tab:related}
	\small
	\begin{tabularx}{\textwidth}{lXccc}
		\hline
		\textbf{Hệ} & \textbf{Phần ngôn ngữ} & \textbf{T.Việt} & \textbf{Thị giác} & \textbf{RL} \\
		\hline
		nanoGPT & tự pretrain (tiếng Anh) & --- & --- & --- \\
		NlpHUST GPT-2 & tự pretrain & \checkmark & --- & --- \\
		PhoGPT-4B & tự pretrain (3,7B, đa GPU) & \checkmark & --- & --- \\
		nanoVLM & SmolLM2 huấn luyện sẵn & --- & \checkmark & --- \\
		LaVy / Vintern-1B & LLM tỷ tham số sẵn & \checkmark & \checkmark & --- \\
		\textbf{ViVLM-nano} & \textbf{tự pretrain (101M, 1 GPU)} & \checkmark & \checkmark & \checkmark \\
		\hline
	\end{tabularx}
\end{table}
